\section{Derivation of the CFE from the Effective Potential}
%=============================================================================

In background-field gauge, the one-loop effective potential for the electromagnetic 
coupling modulus takes the form. Here $U(\alpha)$ is the induced potential for $\sigma$
after rewriting $\alpha=\alpha(\sigma)$:
\begin{equation}
    U(\alpha) = \frac{1}{4}(F_{\mu\nu}F^{\mu\nu})_{\text{eff}} 
        + \text{quantum corrections}
\end{equation}

The full structure, including the Chern--Simons contribution from 11D reduction, is:
\begin{equation}
    U(\alpha) = \frac{A}{4}\alpha^4 - \frac{2}{3}A\cthree^3\alpha^3 
        - 8A\bone\cthree^6 L\,\alpha + \mathcal{O}(\alpha^5)
    \label{eq:potential}
\end{equation}
where $L = \ln(1/\phiz)$ encodes the geometric scale.

The modulus-closure condition $\delta\Gamma_{\mathrm{eff}}/\delta\sigma=0$ reduces to
$\partial U/\partial\alpha = 0$ (monotonic $\alpha(\sigma)$), giving:
\begin{equation}
    \frac{\partial U}{\partial\alpha} = A\alpha^3 - 2A\cthree^3\alpha^2 
        - 8A\bone\cthree^6 L = 0
    \label{eq:stationarity}
\end{equation}

Dividing by $A$ (which is nonzero):
\begin{equation}
    \alpha^3 - 2\cthree^3\alpha^2 - 8\bone\cthree^6 L = 0
    \label{eq:CFE-appendix}
\end{equation}

This is the CFE. The combinatorial factors arise from specific one-loop diagrams:

\textbf{(i) Factor 2 in $2\cthree^3\alpha^2$:}

The cubic term in the effective potential comes from the vacuum polarization 
(``bubble'') diagram with axion insertion:
\begin{center}
\begin{tikzpicture}[scale=0.8]
    \draw[thick, photon] (0,0) -- (1.5,0);
    \draw[thick] (1.5,0) arc (180:0:0.75);
    \draw[thick] (1.5,0) arc (-180:0:0.75);
    \draw[thick, photon] (3,0) -- (4.5,0);
    \draw[thick, dashed] (2.25,0.75) -- (2.25,1.5);
    \node at (2.25,1.8) {\small $a$};
    \node at (0,-0.3) {\small $\gamma$};
    \node at (4.5,-0.3) {\small $\gamma$};
\end{tikzpicture}
\end{center}
The $\varepsilon^{\mu\nu\rho\sigma}\varepsilon_{\alpha\beta\gamma\delta}$ contraction from two 
$F\tilde{F}$ vertices produces a factor 2 from the antisymmetric tensor algebra:
\begin{equation}
    \varepsilon^{\mu\nu\rho\sigma}\varepsilon_{\mu\nu\alpha\beta} = -2(\delta^\rho_\alpha\delta^\sigma_\beta - \delta^\rho_\beta\delta^\sigma_\alpha)
\end{equation}

\textbf{(ii) Factor 8 and $\cthree^6$ in $8\bone\cthree^6 L$:}

The log-dependent term has two components:

\textbf{Origin of $\cthree^6$:} The factor $\cthree^6$ arises from the product structure 
of the effective action. With $A = 2\cthree^3$ as the overall amplitude:
\begin{equation}
    8\bone\cthree^6 L = 8\bone \cdot \cthree^3 \cdot \cthree^3 \cdot L
    = 4 \cdot A \cdot \cthree^3 \cdot \bone L
\end{equation}
The factorization proceeds as follows:
\begin{itemize}[leftmargin=1.5em, topsep=2pt, itemsep=1pt]
    \item One factor of $\cthree^3$ comes from three $\cthree\, aF\tilde{F}$ vertices 
        in the anomaly chain (triangle diagram).
    \item The second factor of $\cthree^3$ comes from the amplitude normalization $A/2 = \cthree^3$.
    \item The coefficient 8 decomposes as $8 = 4 \times 2$, where 4 is from $|D_2|^2$ 
        (two $\mathbb{Z}_2$ reflections) and 2 is from the double cover sheets.
\end{itemize}
\textbf{Consistency check:} $8\cthree^6 = 8 \times (8\pi)^{-6} = 8/(8\pi)^6$, which matches 
the numerical coefficient in the CFE.

\textbf{Origin of factor 8:} The dihedral symmetry $D_4$ of the 4-point amplitude 
contributes $|D_4| = 8$. This arises from the 4 rotations and 4 reflections that 
leave the fermion loop topology invariant:
\begin{center}
\begin{tikzpicture}[scale=0.8]
    \draw[thick] (0,0) rectangle (2,2);
    \draw[thick, photon] (-0.8,0) -- (0,0);
    \draw[thick, photon] (-0.8,2) -- (0,2);
    \draw[thick, photon] (2,0) -- (2.8,0);
    \draw[thick, photon] (2,2) -- (2.8,2);
    \node at (1,-0.5) {\small fermion loop};
    \node at (1,1) {\small $D_4$};
\end{tikzpicture}
\end{center}

\textbf{The $\bone$ factor:} $\bone = 41/10$ counts the hypercharge-weighted fermion 
contributions summed over the SM spectrum. It enters as $\sum_f Y_f^2$, where the 
sum runs over all SM fermions with their multiplicities.

\textbf{(iii) Factor 2/3 in potential:}

The coefficient $2/3$ in the cubic term of \eqref{eq:potential} ensures that 
differentiation produces the factor 2 in \eqref{eq:stationarity}:
\begin{equation}
    \frac{d}{d\alpha}\left(-\frac{2}{3}A\cthree^3\alpha^3\right) = -2A\cthree^3\alpha^2
\end{equation}
This is not a choice but a consequence of the integral representation of the 
effective potential: the cubic term arises from a triangle subgraph.

\textbf{Why no $\alpha^2$ term?}

The absence of an $\alpha^2$ term in the CFE is not accidental. In background-field gauge, 
quadratic terms are absorbed into the field-strength renormalization. The Ward identity 
ensures that physical observables depend only on odd powers of $\alpha$ (gauge invariance 
requires $F^2 \sim \alpha^2$ to enter as a full factor, not as a correction to $\alpha$).

\begin{tcolorbox}[colback=purple!3, colframe=purple!50!black, 
    title={\textbf{Why Integer Coefficients? (No $\zeta(3)$, $\ln 2$, etc.)}}, fonttitle=\bfseries\small]
\textbf{Potential objection:} In typical QFT calculations (e.g., anomalous magnetic moment $g-2$), 
coefficients involve transcendental numbers like $\zeta(3)$, $\pi^2/12$, $\ln 2$. Why are the 
CFE coefficients (2, 8, 48) simple integers?

\textbf{Resolution:} The CFE is not a perturbative loop expansion but a \emph{topological constraint}.

\begin{enumerate}[leftmargin=1.5em, topsep=2pt, itemsep=1pt]
    \item \textbf{Origin of transcendentals:} In $g-2$, transcendentals arise from 
        \emph{integrating over continuous momenta} in Feynman integrals. The integrals 
        produce polylogarithms, which evaluate to $\zeta$-values.
    
    \item \textbf{Topological discreteness:} In TFPT, the relevant integrals are over 
        \emph{discrete topological data}: Euler characteristics, winding numbers, and 
        symmetry group orders. These are integers by definition.
    
    \item \textbf{Combinatorial origin:}
        \begin{itemize}[leftmargin=1em, topsep=1pt, itemsep=1pt]
            \item Factor 2: From $\varepsilon^{\mu\nu\rho\sigma}$ contraction (antisymmetry)
            \item Factor 8: From $|D_4| = 8$ (dihedral group of the box)
            \item Factor 48: \emph{Mnemonic} $2 \times 24$ (double cover $\times$ permutations); the derivation is the spin-lifted deficit with four cut patches (Appendix~\ref{app:48deriv})
        \end{itemize}
    
    \item \textbf{Where $\pi$ appears:} The factors of $\pi$ in $\cthree = 1/(8\pi)$ and 
        $\phitree = 1/(6\pi)$ come from the \emph{normalization} of circles ($2\pi$ per cycle), 
        not from loop integrals. They are geometric, not perturbative.
\end{enumerate}

\textbf{Diagnostic:} If transcendentals appeared in the CFE, it would indicate non-topological 
dynamics. Their absence is evidence that the fixed-point structure is genuinely topological.
\end{tcolorbox}

\begin{remark}
The consistency between \eqref{eq:potential} and \eqref{eq:CFE-appendix} can be 
verified directly: $\frac{d}{d\alpha}\bigl(-\frac{2}{3}A\cthree^3\alpha^3\bigr) 
= -2A\cthree^3\alpha^2$, matching the required coefficient.
\end{remark}

%=============================================================================
\section{Verification Suite Summary}
%=============================================================================

The TFPT verification suite (\texttt{tfpt-suite/}) turns paper claims into deterministic, assumption-explicit computations. All counts in this section are auto-generated from the suite artifacts (no hardcoded ``39 modules''):
\[
\text{unique modules (registry)}=\TFPTSuiteUniqueModulesTotal,\qquad
\text{module artifacts (engineering+physics)}=\TFPTSuiteModulesTotal,\qquad
\text{checks total}=\TFPTSuiteChecksTotal.
\]
It ships:
\begin{itemize}[leftmargin=1.5em, topsep=2pt, itemsep=1pt]
    \item a \textbf{conventional suite} (registry-driven; core invariants, $\alpha$ closure + policy bridge, RG/matching, flavor pipelines, cosmology/bounce wiring, torsion/birefringence, dashboards),
    \item an \textbf{unconventional suite} of search/audit/design tools (for closing explicit ToE gaps without hiding assumptions),
    \item a \textbf{theoryv3} branch of discrete audits that formalizes the ``kernel + crosslinks'' thesis (holonomy $g=5$, defect partition, constant factory).
\end{itemize}
Each module emits \texttt{results.json} (machine-readable), \texttt{report.txt} (human-readable), and optional plots; PDF reports summarize runs.

\subsection{Known Gaps and Status (Auto-Generated)}
\label{sec:known-gaps}

\begin{tcolorbox}[colback=yellow!3, colframe=yellow!60!black,
title={\textbf{Current WARN/FAIL inventory (from suite manifest)}}, fonttitle=\bfseries\small]
\TFPTSuiteProblemModules
\end{tcolorbox}

% Compact table of problem modules (WARN/FAIL/plot issues)
% Inlined (standalone build): formerly \input{../tfpt-suite/out/suite_manifest_table.tex}
% -----------------------------------------------------------------------------
% AUTO-GENERATED by tfpt_suite.suite_manifest.py — do not edit by hand.
\begin{table}[H]
\centering
\small
\begin{tabular}{@{}llrrrrrr@{}}
\toprule
Module & Mode & Checks & PASS & WARN & FAIL & Plots(exp) & Plots(ok)\\
\midrule
\texttt{alpha\_on\_shell\_bridge} & engineering & 7 & 6 & 1 & 0 & 0 & 0\\
\texttt{matching\_finite\_pieces} & engineering & 5 & 4 & 1 & 0 & 0 & 0\\
\texttt{torsion\_condensate} & engineering & 8 & 5 & 3 & 0 & 0 & 0\\
\texttt{alpha\_on\_shell\_bridge} & physics & 7 & 6 & 1 & 0 & 0 & 0\\
\texttt{matching\_finite\_pieces} & physics & 5 & 4 & 1 & 0 & 0 & 0\\
\texttt{torsion\_condensate} & physics & 8 & 5 & 2 & 1 & 0 & 0\\
\bottomrule
\end{tabular}
\caption{Suite manifest: compact table of WARN/FAIL modules and plot issues (auto-generated).}
\end{table}

\subsection{theoryv3 Evidence Layer (Compressed, Auto-Generated)}
\label{sec:theoryv3-evidence}

The \texttt{theoryv3} analysis branch is a compact ``evidence layer'' for the discrete kernel and crosslink logic. Current inventory (auto-generated from \codepath{tfpt-suite/2.5/theoryv3/out/}):
\[
\text{modules}=\TFPTTheoryVThreeModulesTotal,\qquad
\text{checks}=\TFPTTheoryVThreeChecksTotal\ (\mathrm{PASS/WARN/FAIL}=\TFPTTheoryVThreeChecksPassTotal/\TFPTTheoryVThreeChecksWarnTotal/\TFPTTheoryVThreeChecksFailTotal),\qquad
\text{plots expected/present}=\TFPTTheoryVThreePlotsExpectedTotal/\TFPTTheoryVThreePlotsPresentTotal.
\]

% Inlined (standalone build): formerly \input{../tfpt-suite/2.5/theoryv3/out/theoryv3_manifest_table.tex}
% -----------------------------------------------------------------------------
% AUTO-GENERATED by tfpt-suite/2.5/theoryv3/theoryv3_manifest.py — do not edit by hand.
\begin{table}[H]
\centering
\small
\setlength{\tabcolsep}{4pt}
\begin{tabularx}{\linewidth}{@{}>{\raggedright\arraybackslash}p{0.36\linewidth}Xlcc@{}}
\toprule
Module & Signal & Result & Checks & Plots OK\\
\midrule
\texttt{alpha\_backreaction\_sensitivity\_audit} & $\alpha$ backreaction sensitivity (k-grid) & PASS & 2/2 & yes\\
\texttt{axion\_dm\_audit} & Axion DM audit & PASS & 2/2 & yes\\
\texttt{baryon\_consistency\_audit} & Baryon consistency / H0 audit & WARN & 3/4 & yes\\
\texttt{constant\_factory\_audit} & Constant factory audit (derived outputs) & PASS & 3/3 & yes\\
\texttt{dark\_energy\_exponential\_audit} & Dark energy exponential candidate & PASS & 2/2 & yes\\
\texttt{dark\_energy\_norm\_half\_origin\_audit} & Dark energy normalization n=1/2 scan & PASS & 2/2 & yes\\
\texttt{defect\_partition\_g5\_audit} & Two-defect partition ($\delta_2=(5/4)\,\delta_{\rm top}^2$) & PASS & 4/4 & yes\\
\texttt{flavor\_pattern\_audit} & Flavor pattern audit (texture invariants) & PASS & 3/3 & yes\\
\texttt{g5\_crosslink\_audit} & Crosslink consistency (g=5 anchors) & PASS & 3/3 & yes\\
\texttt{g5\_origin\_audit} & Holonomy degeneracy g=5 & PASS & 2/2 & yes\\
\texttt{pmns\_tm1\_audit} & PMNS TM1 audit & PASS & 2/2 & yes\\
\texttt{seed\_invariants\_audit} & Seed invariants (c3, $\varphi_0$, $\delta_{\rm top}$, $\beta$) & PASS & 7/7 & yes\\
\texttt{yukawa\_exponent\_index\_audit} & Yukawa exponent-index mapping & WARN & 3/7 & yes\\
\texttt{yukawa\_index\_mapping\_audit} & Yukawa index mapping & WARN & 3/7 & yes\\
\bottomrule
\end{tabularx}
\caption{theoryv3 compressed evidence layer (auto-generated from \,\codepath{tfpt-suite/2.5/theoryv3/out/}).}
\end{table}

\begin{tcolorbox}[colback=gray!5, colframe=gray!60,
title={\textbf{Two ``lock-in'' results (negative controls; theoryv3)}}, fonttitle=\bfseries\small]
\textbf{$k=2$ is selected by a sensitivity minimum:} the \texttt{alpha\_backreaction\_sensitivity\_audit} shows the minimum $|\mathrm{ppm}|$ in the $k$ grid occurs at $k=2$, with $|\mathrm{ppm}|\simeq 0.037$.

\textbf{$g=5$ is selected by a defect-partition negative control:} the \texttt{defect\_partition\_g5\_audit} shows $g=5$ minimizes $|z|$ sharply versus $g\in\{4,6\}$ (gap $\sim 45$ in $|z|$).
\end{tcolorbox}

\begin{figure}[H]
\centering
\begin{tikzpicture}
\begin{groupplot}[
    group style={group size=2 by 1, horizontal sep=1.2cm},
    width=0.48\textwidth,
    height=0.42\textwidth,
    grid=both,
    legend style={font=\scriptsize},
]
\nextgroupplot[
    title={$k$-sensitivity (abs ppm)},
    xlabel={$k$},
    ylabel={$|{\rm ppm}|$},
    xmin=-0.1, xmax=3.1,
    ymin=0, ymax=4.2,
    xtick={0,1,1.5,2,2.5,3},
]
\addplot+[mark=*, mark size=2.2pt] coordinates {
    (0.0, 3.665371791638453)
    (1.0, 1.8074246412353618)
    (1.5, 0.883514212976573)
    (2.0, 0.03703710120260405)
    (2.5, 0.9542414840493093)
    (3.0, 1.8681110743677323)
};
\addplot[gray, dashed] coordinates {(2.0,0.0) (2.0,4.2)};

\nextgroupplot[
    title={$g$ negative control (z-score)},
    xlabel={$g$},
    ylabel={$z$},
    xmin=3.7, xmax=6.3,
    ymin=-60, ymax=60,
    xtick={4,5,6},
]
\addplot+[mark=*, mark size=2.2pt] coordinates {
    (4, -46.8455434804275)
    (5, 1.8646873709648437)
    (6, 50.5749164827739)
};
\addplot[gray, dashed] coordinates {(5,-60) (5,60)};
\end{groupplot}
\end{tikzpicture}
\caption{theoryv3 lock-in diagnostics: $k$-sensitivity of $\alpha^{-1}(0)$ (left) and the $g\in\{4,5,6\}$ negative control selecting $g=5$ (right).}
\end{figure}

\begin{tcolorbox}[colback=gray!5, colframe=gray!60,
title={\textbf{Global scorecard (dashboard $\chi^2$; no covariance)}}, fonttitle=\bfseries\small]
The \texttt{global\_consistency\_test} module evaluates a reference-table $\chi^2$ over a small set of headline observables (engineering mode). Current totals:
\[
\chi^2_{\mathrm{core}}=3.083\ (6~\mathrm{dof}),\qquad p\simeq 0.798,
\]
and (excluding $\alpha$ metrology dominance)
\[
\chi^2_{\mathrm{core}\setminus\alpha}=1.354\ (5~\mathrm{dof}),\qquad p\simeq 0.929.
\]
This is a dashboard (independent Gaussians; simplified $r$ bound proxy), not a publication-grade global likelihood. The suite reports per-term contributions and an ``excluding $\alpha$'' total explicitly to avoid misleading optics when a single ultra-precise reference dominates.
\end{tcolorbox}

\begin{figure}[H]
\centering
\begin{tikzpicture}
\begin{groupplot}[
    group style={group size=2 by 1, horizontal sep=1.2cm},
    ybar,
    bar width=10pt,
    width=0.48\textwidth,
    height=0.42\textwidth,
    xmin=0.5, xmax=4.5,
    xtick={1,2,3,4},
    xticklabels={$\beta_{\rm deg}$, $\lambda$, $n_s$, $A_s$},
    x tick label style={rotate=0, font=\scriptsize},
    grid=both,
]
\nextgroupplot[
    title={$\chi^2$ contributions (excluding $\alpha$)},
    ylabel={$\chi^2$},
    ymin=0, ymax=0.8,
]
\addplot coordinates {
    (1,0.5903174784328564)
    (2,0.6542656358088781)
    (3,0.021391549817205795)
    (4,0.08789723472938661)
};

\nextgroupplot[
    title={$z$-scores (excluding $\alpha$)},
    ylabel={$z$},
    ymin=-1.0, ymax=0.15,
]
\addplot coordinates {
    (1,-0.7683212078505034)
    (2,-0.8088668838621581)
    (3,-0.14625850340136054)
    (4,-0.296474678057649)
};
\addplot[black!50, dashed] coordinates {(0.5,0) (4.5,0)};
\end{groupplot}
\end{tikzpicture}
\caption{Engineering-mode dashboard plots from \texttt{global\_consistency\_test}: $\chi^2$ contributions (left) and $z$-scores (right), both excluding $\alpha$.}
\end{figure}

%=============================================================================
\section{\texorpdfstring{E$_8$ Chain Certificate}{E8 Chain Certificate}}
\label{app:E8cert}
%=============================================================================

This appendix provides the complete verification data for the E$_8$ chain uniqueness 
theorem (Theorem~\ref{thm:chainunique}).

\subsection{Orbit Data Structure}

Each nilpotent orbit $\mathcal{O}$ of $\mathfrak{e}_8$ is characterized by:
\begin{itemize}[leftmargin=1.5em, topsep=2pt, itemsep=1pt]
    \item \textbf{Bala--Carter label:} Standard notation (e.g., $A_1$, $2A_1$, $A_2$, \ldots)
    \item \textbf{Dimension:} $\dim(\mathcal{O})$
    \item \textbf{Centralizer dimension:} $D(\mathcal{O}) = 248 - \dim(\mathcal{O})$
    \item \textbf{Height:} $h(\mathcal{O})$ (sum of weighted Dynkin diagram coefficients)
\end{itemize}

